\section{Introduction}

The sum of the linear Fourier coefficients of a Boolean function measures
its correlation with the sum of its inputs. For a real-valued function
\(f:\{-1,1\}^n\to\mathbb R\), with \(X\) uniform on the sign cube,
write
\[
 \widehat f(S)=\E\!\left[f(X)\prod_{i\in S}X_i\right],
 \qquad \deg(f)=\max\{|S|:\widehat f(S)\ne0\}.
\]
Constants have degree zero. This is the degree of the real multilinear
polynomial representing \(f\), as in the work of Nisan and
Szegedy~\cite{NisanSzegedy}. For Boolean \(f\), taking values in
\(\{-1,1\}\), Cauchy--Schwarz gives
\[
 \sum_{i=1}^n\widehat f(\{i\})
 =\E\!\left[f(X)\sum_{i=1}^nX_i\right]\le\sqrt n.
\]
Gopalan and Servedio conjectured that \(\sqrt{\deg(f)}\) could replace
\(\sqrt n\); see O'Donnell's problem list~\cite[p.~9]{ODonnell}.
We show that no constant multiple of this proposed bound holds.

\begin{theorem}\label{thm:main}
For every real \(C>0\), there are a positive integer \(n\) and a
nonconstant Boolean function \(f:\{-1,1\}^n\to\{-1,1\}\) such that
\[
 \sum_{i=1}^n\widehat f(\{i\})>C\sqrt{\deg(f)}.
\]
Consequently,
\[
 \sup_{\substack{n\ge1,\ f:\{-1,1\}^n\to\{-1,1\}\\\deg(f)>0}}
 \frac{\sum_i|\widehat f(\{i\})|}{\sqrt{\deg(f)}}=\infty.
\]
\end{theorem}

The universal signed and absolute versions of the conjecture are
equivalent. Indeed, the coordinate change
\(x_i\mapsto\operatorname{sign}(\widehat f(\{i\}))x_i\), with sign
\(+1\) at zero, preserves degree and makes every linear coefficient
nonnegative. Our construction produces a large signed sum directly.
The dimension and all copy counts in the
construction are finite; we do not obtain useful bounds on their growth.

\subsection*{The square-root and majority bounds}
The conjecture is attributed to Gopalan and Servedio, circa 2009, in
O'Donnell's problem list. It also appears as Conjecture 3.17 in the
preprint of Filmus, Hatami, Keller and Lifshitz~\cite{FHKL}.
Here the singleton coefficients are summed linearly, rather than squared.

A sharper proposal in the same problem list uses the majority benchmark
\[
 B_d=\E|X_1+\cdots+X_d|\le\sqrt d.
\]
This is the signed singleton sum of majority on \(d\) bits, with either
fixed value at ties. Replacing \(\sqrt{\deg(f)}\) by
\(B_{\deg(f)}\) gives a stronger inequality, so a counterexample to
that inequality need not violate the square-root bound.
Jha~\cite[Theorem 2.1]{Jha} gave discrete-derivative reformulations of
the majority proposal. Wang~\cite[Theorems 2.2, 3.2 and 3.6]{Wang}
gave another equivalent formulation and proved the cases \(d=1\) and
\(d=n-1\). Kudin and Pa\v{s}ali\'c~\cite{KudinPasalic} proved the
cases \(d=2,3\) and refuted the majority bound at \(d=4\).
Theorem~\ref{thm:main} excludes every constant multiple of the weaker
square-root bound. Stating both inequalities explicitly distinguishes
the questions despite the different names used for them in the literature.

A linear degree bound does hold for every Boolean function. Let
\(\operatorname{Inf}_i(f)\) be the probability that flipping coordinate
\(i\) changes its value, and let
\(I(f)=\sum_i\operatorname{Inf}_i(f)\) be the total influence.
The discrete derivative in coordinate \(i\), defined as half the
difference between the values with that coordinate set to \(+1\) and
\(-1\), takes values in \(\{-1,0,1\}\). Its mean is
\(\widehat f(\{i\})\), so the absolute value of this mean is at most
its second moment, \(\operatorname{Inf}_i(f)\). The Fourier formula
for total influence and Parseval's identity therefore give
\cite[Lemma~2.4 and Corollary~2.5]{NisanSzegedy}
\[
 \sum_i|\widehat f(\{i\})|
 \le I(f)=\sum_{S\subseteq[n]}|S|\widehat f(S)^2
 \le\deg(f).
\]
Theorem~\ref{thm:main} thus separates the conjectured square-root
scale from this surviving linear bound.

\subsection*{Retaining variance with low-degree observations}
Instead of constructing \(f\) immediately, we first construct a
finite-valued observation \(F\) of independent signs. If
\(H_N=X_1+\cdots+X_N\), the information that \(F\) retains about this
sum is measured by
\[
 T_F=\E[H_N\mid F],\qquad v(F)=\E T_F^2.
\]
We call \(T_F\) the \emph{score} of \(F\) and \(v(F)\) its retained
variance. When \(v(F)>0\), its standardized score is
\(T_F/\sqrt{v(F)}\). We call a positive real number \(D\) a \emph{cell degree bound} if
 every output indicator \(\one_{\{F=a\}}\) has degree at most \(D\). Every function of
\(F\) then has degree at most \(D\), including
\(f=\operatorname{sign}(T_F)\). This readout has correlation
\(\E[fH_N]=\E|T_F|\). Averaging independent copies makes its score
approximately Gaussian, so a large ratio \(v(F)/D\) gives the theorem.

Conditional coordinate sums also appear in decision-tree inequalities.
A subcube is obtained by fixing some coordinates. If the cells of
\(F\) are subcubes, the unfixed coordinates remain independent uniform
signs under conditioning, so \(v(F)\) equals the expected number of
fixed coordinates. O'Donnell and Servedio~\cite[Theorem~3 and
Lemma~3]{ODonnellServedio} bound the signed singleton sum of a Boolean
readout by the square root of this expectation. Blais, Tan and
Wan~\cite[Lemma~3.1 and Theorem~1]{BlaisTanWan} use the conditional
coordinate sum to prove square-root bounds in parity decision-tree
depth, where a query reveals the product of a specified set of inputs.
Our cells need not be subcubes: their indicators are controlled by
polynomial degree, and cancellation allows retained variance to exceed
that degree cost.

The main construction repeatedly increases this ratio. Revealing
\(m\) independent standardized scores retains variance \(m\). We use
\(m+1\) scores and design an observation retaining strictly more than
\(m\) units of variance, although each output indicator is a linear
combination of functions of at most \(m\) inputs. Such an indicator may
depend on all \(m+1\) inputs: the degree saving comes from cancellation,
not from ignoring a fixed coordinate.

Here is the reporting rule behind the gain. One score selects a leaf
among the remaining \(m\) scores, each of which has a short prescribed
interval. Usually we report all but one score. On the exceptional event
that only the selected leaf misses its interval, we report a single
tag. The intervals are chosen so that the sum of all the scores can be
predicted on this event with mean squared error less than one. Every
other report has residual error exactly one. Thus the exceptional tag
improves on the variance retained by revealing \(m\) scores.
Its indicator admits the required cancellation identity.

Section~\ref{sec:observations} proves the finite-observation calculus
and amplification criterion. Section~\ref{sec:observation} then
explains the reports and their
prediction error before choosing the Gaussian intervals. Once that
choice is fixed, the central limit theorem supplies a sufficiently large
finite average at each step. A finite number of steps gives any
prescribed ratio, and a final sign readout proves the theorem. Thus
Sections~\ref{sec:observations} and~\ref{sec:observation} contain a
complete proof of Theorem~\ref{thm:main}.

The further constructions distinguish which features of this rule are
needed for a gain. In Section~\ref{sec:alternatives}, the \emph{coarse
report} shows that the all-leaves-pass values can also be replaced by
one tag; switching to an extra independent score isolates the
conditional-variance saving on the exceptional event.
Section~\ref{sec:mean-cost} allows different degree costs for different
outputs: retained variance need only exceed their \emph{average} cost.
Reporting typical words of outputs and merging the remaining words into
one cell then gives a degree bound valid everywhere while retaining
almost all the variance.
Appendix~\ref{app:weighted-mixture} supplies an independent Gaussian
event for the average-cost criterion. Appendix~\ref{app:small-width}
gives a different Gaussian parameter schedule for the coarse report,
together with the transfer argument for its law-dependent coefficients.
